\section{Locality and Degrees of Freedom}

Most classical field theories are local. This means that the equation at a point
depends on the value of the field and finitely many derivatives of the field at
that point. For a scalar field, a typical local equation may involve
\[
    \phi, \qquad
    \frac{\partial \phi}{\partial t}, \qquad
    \nabla \phi, \qquad
    \frac{\partial^2 \phi}{\partial t^2}, \qquad
    \Delta \phi.
\]
It does not directly require the value of \(\phi\) at every distant point.

This locality principle is one reason differential equations are central in
field theory. Partial derivatives compare nearby field values, so local
physical laws become partial differential equations.

\begin{summarybox}
A finite-dimensional mechanical system has coordinates \(q^i(t)\). A field has
degrees of freedom labelled by position as well as by time, such as
\(\phi(t,\vect{x})\). Classical field theory is therefore the natural language
for waves, fluids, electromagnetic fields, and relativistic continuous systems.
\end{summarybox}
